\section{The key technical lemma}
\label{sec:lemma1}

This section proves the single-bit conditional-entropy gap that drives
the main result. The setup: $C$ is a random variable on a finite set $S$,
and for each $c \in S$ we have a parameter $p_c \in [0, 1]$. Conditional on
$C = c$, sample a Bernoulli $X$ with $\Pr[X = 1 \mid C = c] = p_c$. Let
$(C', X')$ be an independent copy of $(C, X)$, so $X'$ is conditionally
Bernoulli$(p_{C'})$ given $C'$ and is independent of $(C, X)$.

In the application to the union-closed conjecture (\Cref{sec:theorem1}),
$C$ will be the history $A_{<i}$ of bits revealed so far, $p_c$ will be the
conditional probability $\Pr[A_i = 1 \mid A_{<i} = c]$, and the
expectation hypothesis $\E[X] \le 0.01$ will correspond to $\Pr[i \in A] \le 0.01$.

\begin{lemma}[Key conditional-entropy gap]\label{lem:gap}
With the setup above, assume $\E[X] \le 0.01$. Then
\begin{align}\label{eq:lem-gap}
\H(X \cup X' \mid C, C') \;\ge\; 1.26 \, \H(X \mid C).
\end{align}
Equivalently, writing $q(c) := \Pr[C = c]$ and using that
$X \cup X' \mid (C, C') = (c, c')$ is $\Bern(p_c + p_{c'} - p_c p_{c'})$,
\begin{align}\label{eq:lem-gap-explicit}
\sum_{c, c' \in S} q(c) q(c') \, \hbin(p_c + p_{c'} - p_c p_{c'})
\;\ge\; 1.26 \, \sum_{c \in S} q(c) \, \hbin(p_c).
\end{align}
\end{lemma}

\begin{proof}
The proof partitions the values of $c$ into two regimes by whether $p_c$ is
small or large, then applies \Cref{lem:hbin-low} to the both-small block and
\Cref{lem:hbin-mixed} to the mixed (one-small, one-large) block, and finally
discards the both-large block (whose contribution is non-negative).

\paragraph{Step 1: Partition by $p_c$ and Markov.}
Define
\begin{align*}
\Ccal_0 \;:=\; \{ c \in S : p_c \le 0.1 \}, \qquad
\Ccal_1 \;:=\; S \setminus \Ccal_0 \;=\; \{ c \in S : p_c > 0.1 \}.
\end{align*}
Since $\E[X] = \sum_c q(c) p_c \le 0.01$, Markov's inequality gives
\begin{align}\label{eq:markov}
\Pr[C \in \Ccal_1] \;=\; \sum_{c : p_c > 0.1} q(c)
\;\le\; \frac{\E[X]}{0.1} \;\le\; \frac{0.01}{0.1} \;=\; 0.1,
\end{align}
and therefore
\begin{align}\label{eq:c0-mass}
\Pr[C \in \Ccal_0] \;\ge\; 0.9.
\end{align}

We abbreviate $\Pr[\Ccal_0] := \Pr[C \in \Ccal_0]$,
$\Pr[\Ccal_0'] := \Pr[C' \in \Ccal_0]$,
$\Pr[\Ccal_0, \Ccal_0'] := \Pr[C \in \Ccal_0 \wedge C' \in \Ccal_0]$,
and similarly for $\Ccal_1$. By independence of $C$ and $C'$ and the fact
that $C'$ has the same distribution as $C$,
$\Pr[\Ccal_0] = \Pr[\Ccal_0']$,
$\Pr[\Ccal_0, \Ccal_0'] = \Pr[\Ccal_0]^2$,
$\Pr[\Ccal_0, \Ccal_1'] = \Pr[\Ccal_0] \Pr[\Ccal_1]$,
$\Pr[\Ccal_1, \Ccal_0'] = \Pr[\Ccal_1] \Pr[\Ccal_0]$,
and $\Pr[\Ccal_1, \Ccal_1'] = \Pr[\Ccal_1]^2$.

\paragraph{Step 2: Decompose both sides of \cref{eq:lem-gap-explicit}.}
The right-hand side splits as
\begin{align}\label{eq:rhs-decomp}
\H(X \mid C)
\;=\; \sum_{c \in S} q(c) \, \hbin(p_c)
\;=\; \Pr[\Ccal_0] \, \H(X \mid C \in \Ccal_0) \;+\; \Pr[\Ccal_1] \, \H(X \mid C \in \Ccal_1),
\end{align}
where $\H(X \mid C \in \Ccal_j) := \sum_{c \in \Ccal_j} \frac{q(c)}{\Pr[\Ccal_j]} \hbin(p_c)$.
The left-hand side splits by the partition of $(C, C')$ into the four
events $\{\Ccal_0 \cap \Ccal_0',\, \Ccal_0 \cap \Ccal_1',\, \Ccal_1 \cap \Ccal_0',\, \Ccal_1 \cap \Ccal_1'\}$:
\begin{align}
\H(X \cup X' \mid C, C')
&\;=\; \Pr[\Ccal_0, \Ccal_0'] \, \H(X \cup X' \mid \Ccal_0, \Ccal_0')
\;+\; \Pr[\Ccal_0, \Ccal_1'] \, \H(X \cup X' \mid \Ccal_0, \Ccal_1') \notag\\
&\quad +\; \Pr[\Ccal_1, \Ccal_0'] \, \H(X \cup X' \mid \Ccal_1, \Ccal_0')
\;+\; \Pr[\Ccal_1, \Ccal_1'] \, \H(X \cup X' \mid \Ccal_1, \Ccal_1'). \label{eq:lhs-decomp}
\end{align}
By symmetry under swapping $(C, X) \leftrightarrow (C', X')$ (which leaves
$X \cup X' = \max(X, X')$ invariant),
$\Pr[\Ccal_0, \Ccal_1'] \H(X \cup X' \mid \Ccal_0, \Ccal_1')
= \Pr[\Ccal_1, \Ccal_0'] \H(X \cup X' \mid \Ccal_1, \Ccal_0')$, so we can
write the mixed contribution as $2 \Pr[\Ccal_0, \Ccal_1'] \, \H(X \cup X' \mid \Ccal_0, \Ccal_1')$.

\paragraph{Step 3: Bound the both-small block (Claim A).}
We claim
\begin{align}\label{eq:both-small}
\Pr[\Ccal_0]^2 \, \H(X \cup X' \mid \Ccal_0, \Ccal_0')
\;\ge\; 1.26 \, \Pr[\Ccal_0] \, \H(X \mid C \in \Ccal_0).
\end{align}
Let $q_0(c) := q(c) / \Pr[\Ccal_0]$ for $c \in \Ccal_0$ be the conditional
distribution of $C$ given $C \in \Ccal_0$. Then
\begin{align}
\Pr[\Ccal_0] \, \H(X \mid C \in \Ccal_0)
&\;=\; \Pr[\Ccal_0] \, \E_{c \sim q_0}[\hbin(p_c)] \notag \\
&\;=\; \Pr[\Ccal_0] \, \E_{c, c' \sim q_0} \!\left[ \tfrac{\hbin(p_c) + \hbin(p_{c'})}{2} \right] \label{eq:sym-avg}\\
&\;\le\; \frac{\Pr[\Ccal_0]}{1.4} \, \E_{c, c' \sim q_0} \!\left[ \hbin(p_c + p_{c'} - p_c p_{c'}) \right] \label{eq:apply-lemma-low}\\
&\;\le\; \frac{\Pr[\Ccal_0]^2}{0.9 \cdot 1.4} \, \E_{c, c' \sim q_0} \!\left[ \hbin(p_c + p_{c'} - p_c p_{c'}) \right] \label{eq:multiply-c0}\\
&\;=\; \frac{\Pr[\Ccal_0]^2}{1.26} \, \H(X \cup X' \mid \Ccal_0, \Ccal_0'), \label{eq:identify-cup}
\end{align}
where (i) Eq.~\eqref{eq:sym-avg} is the symmetrization
$\E_c[\hbin(p_c)] = \E_c[\hbin(p_c)/2] + \E_{c'}[\hbin(p_{c'})/2]
= \E_{c,c'}[(\hbin(p_c) + \hbin(p_{c'}))/2]$ valid since $c, c' \sim q_0$ are iid;
(ii) Eq.~\eqref{eq:apply-lemma-low} applies \Cref{lem:hbin-low} pointwise
(since $p_c, p_{c'} \le 0.1$ for $c, c' \in \Ccal_0$) and integrates;
(iii) Eq.~\eqref{eq:multiply-c0} uses
$\Pr[\Ccal_0] \le \Pr[\Ccal_0]^2 / 0.9$, equivalent to
$\Pr[\Ccal_0] \ge 0.9$, which is Eq.~\eqref{eq:c0-mass};
(iv) Eq.~\eqref{eq:identify-cup} identifies
$\E_{c, c' \sim q_0}[\hbin(p_c + p_{c'} - p_c p_{c'})] = \H(X \cup X' \mid \Ccal_0, \Ccal_0')$
since $X \cup X' \mid (C, C') = (c, c')$ is $\Bern(p_c + p_{c'} - p_c p_{c'})$.
Rearranging Eq.~\eqref{eq:identify-cup} gives Eq.~\eqref{eq:both-small}.

\paragraph{Step 4: Bound the mixed block (Claim B).}
We claim
\begin{align}\label{eq:mixed}
2 \, \Pr[\Ccal_0, \Ccal_1'] \, \H(X \cup X' \mid \Ccal_0, \Ccal_1')
\;\ge\; 1.62 \, \Pr[\Ccal_1] \, \H(X \mid C \in \Ccal_1).
\end{align}
By the Bernoulli identity, $\H(X \cup X' \mid (C, C') = (c, c'))
= \hbin(p_c + p_{c'} - p_c p_{c'})$, so
\begin{align}
2 \, \Pr[\Ccal_0, \Ccal_1'] \, \H(X \cup X' \mid \Ccal_0, \Ccal_1')
&\;=\; 2 \sum_{c \in \Ccal_0,\, c' \in \Ccal_1} q(c) q(c') \, \hbin(p_c + p_{c'} - p_c p_{c'}) \notag\\
&\;\ge\; 2 \sum_{c \in \Ccal_0,\, c' \in \Ccal_1} q(c) q(c') \, (1 - p_c) \, \hbin(p_{c'}) \label{eq:apply-mixed}\\
&\;=\; 2 \left( \sum_{c' \in \Ccal_1} q(c') \, \hbin(p_{c'}) \right) \left( \sum_{c \in \Ccal_0} q(c) (1 - p_c) \right) \notag\\
&\;=\; 2 \, \Pr[\Ccal_1] \, \H(X \mid C \in \Ccal_1) \cdot \sum_{c \in \Ccal_0} q(c) (1 - p_c). \label{eq:factor-mixed}
\end{align}
Now for $c \in \Ccal_0$, $p_c \le 0.1$, so $1 - p_c \ge 0.9$, and therefore
\begin{align}\label{eq:sum-1-pc}
\sum_{c \in \Ccal_0} q(c) (1 - p_c) \;\ge\; 0.9 \sum_{c \in \Ccal_0} q(c) \;=\; 0.9 \, \Pr[\Ccal_0]
\;\ge\; 0.9 \cdot 0.9 \;=\; 0.81.
\end{align}
Substituting Eq.~\eqref{eq:sum-1-pc} into Eq.~\eqref{eq:factor-mixed},
\begin{align*}
2 \, \Pr[\Ccal_0, \Ccal_1'] \, \H(X \cup X' \mid \Ccal_0, \Ccal_1')
&\;\ge\; 2 \, \Pr[\Ccal_1] \, \H(X \mid C \in \Ccal_1) \cdot 0.81 \\
&\;=\; 1.62 \, \Pr[\Ccal_1] \, \H(X \mid C \in \Ccal_1),
\end{align*}
which is Eq.~\eqref{eq:mixed}. The inequality in Eq.~\eqref{eq:apply-mixed}
applies \Cref{lem:hbin-mixed} pointwise: for $(c, c') \in \Ccal_0 \times \Ccal_1$,
$\hbin(p_c + p_{c'} - p_c p_{c'}) \ge (1 - p_c) \hbin(p_{c'})$ with the roles
$(p, p') \leftarrow (p_c, p_{c'})$.

\paragraph{Step 5: Assemble the four blocks.}
Using the decomposition in Eq.~\eqref{eq:lhs-decomp} together with the
non-negativity of the both-large block $\Pr[\Ccal_1]^2 \H(X \cup X' \mid \Ccal_1, \Ccal_1') \ge 0$,
\begin{align*}
\H(X \cup X' \mid C, C')
&\;\ge\; \Pr[\Ccal_0]^2 \, \H(X \cup X' \mid \Ccal_0, \Ccal_0')
\;+\; 2 \, \Pr[\Ccal_0, \Ccal_1'] \, \H(X \cup X' \mid \Ccal_0, \Ccal_1') \\
&\;\ge\; 1.26 \, \Pr[\Ccal_0] \, \H(X \mid C \in \Ccal_0)
\;+\; 1.62 \, \Pr[\Ccal_1] \, \H(X \mid C \in \Ccal_1) \\
&\;\ge\; 1.26 \, \Pr[\Ccal_0] \, \H(X \mid C \in \Ccal_0)
\;+\; 1.26 \, \Pr[\Ccal_1] \, \H(X \mid C \in \Ccal_1) \\
&\;=\; 1.26 \, \H(X \mid C),
\end{align*}
where the first inequality drops the non-negative both-large block, the
second applies Claim A (\Cref{eq:both-small}) to the first summand and
Claim B (\Cref{eq:mixed}) to the second, the third uses $1.62 \ge 1.26$ to
weaken the second summand (so that we may pull out a common factor of
$1.26$), and the final equality is Eq.~\eqref{eq:rhs-decomp}.
This establishes Eq.~\eqref{eq:lem-gap} and completes the proof.
\end{proof}
